% !TEX root = ../main.tex
\corrigesChapitre{20}{Transferts thermiques et bilans}
\corrigesSouples

\begin{blocLivre}{cor-20-1}
\begin{corrige}{1}{Trois modes de transfert}
\demarche{Conduction : sans déplacement de matière ; convection : avec mouvement d'un fluide ; rayonnement : ondes
électromagnétiques.}
\begin{questions}
\item La \res{\text{conduction}}, dans le métal, de proche en proche.
\item L'air chauffé au contact du radiateur s'élève et circule dans toute la pièce : c'est la \res{\text{convection}}.
\item Par \res{\text{rayonnement}}, le seul mode de transfert possible dans le vide.
\item Le vide entre les parois supprime la \res{\text{conduction}} et la \res{\text{convection}} ; l'argenture réfléchit le
rayonnement et limite les transferts par \res{\text{rayonnement}} ; le bouchon limite la convection par l'ouverture.
\end{questions}
\end{corrige}
\end{blocLivre}
\begin{blocLivre}{cor-20-2}
\begin{corrige}{2}{La porte d'entrée}
\demarche{$R_\mathrm{th}=e/(\lambda\,S)$ puis $\Phi=\Delta T/R_\mathrm{th}$ et $Q=\Phi\,\Delta t$.}
\begin{questions}
\item $R_\mathrm{th}=\dfrac{0{,}040}{0{,}15\times2{,}0}=\res{\qty{0,133}{K.W^{-1}}}$.
\item $\Phi=\dfrac{20-1}{0{,}133}=\res{\qty{143}{W}}$, de l'intérieur, plus chaud, \res{\text{vers l'extérieur}}.
\item $Q=143\times86\,400=\res{\qty{12,3}{MJ}}$, soit \res{\qty{3,42}{kWh}}.
\item $R_\mathrm{th}$ est proportionnelle à $e$ : il faut doubler l'épaisseur, soit \res{\qty{8,0}{cm}}.
\end{questions}
\end{corrige}
\end{blocLivre}
\begin{blocLivre}{cor-20-3}
\begin{corrige}{3}{Rayonner selon sa température}
\demarche{$P/S=\sigma\,T^4$, avec $T$ en kelvins.}
\begin{questions}
\item $\dfrac{P}{S}=\num{5.67e-8}\times5772^4=\res{\qty{6,29e7}{W.m^{-2}}}$.
\item Sol : \res{\qty{390}{W.m^{-2}}} ; peau : \res{\qty{497}{W.m^{-2}}}.
\item $\left(\dfrac{5772}{288}\right)^4=\res{\num{1.61e5}}$ fois plus.
\item Par $2^4=\res{16}$ : la puissance rayonnée croît très vite avec la température.
\end{questions}
\end{corrige}
\end{blocLivre}
\begin{blocLivre}{cor-20-4}
\begin{corrige}{4}{Le thé refroidit}
\demarche{À $t=\tau$ : $T=T_\mathrm{ext}+0{,}37\,(T_0-T_\mathrm{ext})$ ; la tangente à l'origine coupe l'asymptote en $t=\tau$.}
\begin{questions}
\item $T_0=\res{\qty{85}{\degreeCelsius}}$ ; $T$ tend vers \res{\qty{20}{\degreeCelsius}}, la température de la pièce.
\item La tangente à l'origine coupe la droite $T=\qty{20}{\degreeCelsius}$ en $t=\qty{20}{min}$ ; et $T=20+0{,}37\times65=\qty{43,9}{\degreeCelsius}$
est atteinte à $t=\qty{20}{min}$ : $\res{\tau=\qty{20}{min}}$.
\item $T(2\tau)=20+65\,\mathrm{e}^{-2}=\res{\qty{28,8}{\degreeCelsius}}$ et $T(3\tau)=\res{\qty{23,2}{\degreeCelsius}}$.
\item $60=20+65\,\mathrm{e}^{-t/20}$ donne $t=20\ln\dfrac{65}{40}=\res{\qty{9,7}{min}}$.
\end{questions}
\schema[corrige]{20-ex04}
\end{corrige}
\end{blocLivre}
\begin{blocLivre}{cor-20-5}
\begin{corrige}{5}{L'albédo des surfaces}
\demarche{Puissance absorbée par mètre carré : $(1-A)\times\qty{340}{W.m^{-2}}$.}
\begin{questions}
\item L'albédo est la \res{\text{fraction du rayonnement reçu qui est réfléchie}} par la surface.
\item Neige : \res{\qty{51}{W.m^{-2}}} ; forêt : \res{\qty{289}{W.m^{-2}}} ; océan : \res{\qty{320}{W.m^{-2}}}.
\item L'océan s'échauffe le plus, la neige \res{\text{le moins}}.
\item L'océan absorbe six fois plus que la neige : la fonte de la glace augmente l'énergie absorbée, ce qui accélère la
fonte. C'est une \res{\text{rétroaction positive}}.
\end{questions}
\end{corrige}
\end{blocLivre}
\begin{blocLivre}{cor-20-6}
\begin{corrige}{6}{Simple ou double vitrage}
\demarche{$\Phi=\Delta T/R_\mathrm{th}$ ; énergie $=\Phi\times\Delta t$.}
\begin{questions}
\item Simple : \res{\qty{174}{W}} ; double : \res{\qty{84}{W}} ; double performant : \res{\qty{33}{W}}.
\item $U=\dfrac{1}{0{,}115\times1{,}5}=\res{\qty{5,8}{W.m^{-2}.K^{-1}}}$ : chaque mètre carré perd \qty{5,8}{W} par kelvin d'écart.
\item $\Delta\Phi=\dfrac{12}{0{,}115}-\dfrac{12}{0{,}606}=\qty{84,5}{W}$, d'où $E=84{,}5\times3600=\res{\qty{304}{kWh}}$.
\item $304\times0{,}20\approx\res{61}$~€ par fenêtre et par hiver.
\end{questions}
\end{corrige}
\end{blocLivre}
\begin{blocLivre}{cor-20-7}
\begin{corrige}{7}{Isoler les combles}
\demarche{Couches en série : $R=R_\mathrm{plafond}+R_\mathrm{laine}$.}
\begin{questions}
\item $R_\mathrm{laine}=\dfrac{0{,}30}{0{,}040\times80}=\res{\qty{0,094}{K.W^{-1}}}$.
\item Avant : $\dfrac{20}{\num{3.0e-3}}=\res{\qty{6,67e3}{W}}$ ; après : $\dfrac{20}{\num{3.0e-3}+\num{0.094}}=\res{\qty{207}{W}}$, trente fois moins.
\item $\dfrac{e}{\lambda}=\dfrac{0{,}30}{0{,}040}=\qty{7,5}{m^2.K.W^{-1}}$ : \res{\text{éligible}}. Épaisseur minimale :
$7\times0{,}040=\res{\qty{28}{cm}}$.
\item $\Delta\Phi=\dfrac{12}{\num{3.0e-3}}-\dfrac{12}{\num{0.0968}}=\qty{3,88e3}{W}$, soit $E=\res{\qty{1,40e4}{kWh}}$ par hiver.
\end{questions}
\end{corrige}
\end{blocLivre}
\begin{blocLivre}{cor-20-8}
\begin{corrige}{8}{Modéliser un refroidissement en Python}
\demarche{Méthode d'Euler : $T(t+\mathrm{d}t)\approx T(t)+\dfrac{\mathrm{d}T}{\mathrm{d}t}\,\mathrm{d}t$.}
\begin{questions}
\item Elle traduit $T(t+\mathrm{d}t)=T(t)+\dfrac{T_\mathrm{ext}-T}{\tau}\,\mathrm{d}t$, c'est-à-dire \res{\text{l'équation différentielle}}
du refroidissement, obtenue par le premier principe et la loi de Newton.
\item La méthode d'Euler suppose la dérivée constante pendant chaque pas : c'est une \res{\text{approximation}}. On réduit
l'écart en diminuant \code{dt}.
\item Remplacer la condition par \code{while T > 30:} et afficher \code{t} : environ \res{\qty{2,69e3}{s}}. Par le calcul :
$t=\tau\ln\dfrac{80-20}{30-20}=1500\ln6=\qty{2,69e3}{s}$, soit \qty{44,8}{min}.
\item Créer les listes \code{temps, temperatures = [0], [T0]}, y ajouter \code{t / 60} et \code{T} dans la boucle,
puis \code{plt.plot(temps, temperatures)} et \code{plt.show()}.
\end{questions}
\schema[corrige]{rev-20-ex08-cor}
\legende{Le pas de 10 s produit une courbe très proche de la solution exacte ; les deux tracés se superposent presque.}
\end{corrige}
\end{blocLivre}
\begin{blocLivre}{cor-20-9}
\begin{corrige}{9}{La température de la Terre sans atmosphère}
\demarche{Équilibre radiatif : $(1-A)\,S_0\,\pi R^2=\sigma\,T^4\times4\pi R^2$.}
\begin{questions}
\item Puissance absorbée : $(1-A)\,S_0\,\pi R^2$ ; puissance rayonnée : $\sigma\,T^4\times4\pi R^2$.
\item $T=\left(\dfrac{0{,}70\times1361}{4\times\num{5.67e-8}}\right)^{1/4}=\res{\qty{255}{K}}$, soit environ \res{\qty{-18}{\degreeCelsius}}.
\item Terre glacée : \res{\qty{221}{K}} ; Terre noire : \res{\qty{278}{K}}.
\item L'\res{\text{effet de serre}} : l'atmosphère renvoie vers le sol une partie du rayonnement infrarouge qu'il émet.
\end{questions}
\end{corrige}
\end{blocLivre}
\begin{blocLivre}{cor-20-10}
\begin{corrige}{10}{Froid près d'une fenêtre}
\demarche{Puissance nette perdue : $\sigma\,S\,(T_\mathrm{peau}^4-T_\mathrm{paroi}^4)$, températures en kelvins.}
\begin{questions}
\item $P=\num{5.67e-8}\times1{,}8\times306^4=\res{\qty{895}{W}}$.
\item $P_\mathrm{nette}=\num{5.67e-8}\times1{,}8\times(306^4-293^4)=\res{\qty{143}{W}}$.
\item Face à une paroi à \qty{283}{K} : $P_\mathrm{nette}=\res{\qty{240}{W}}$.
\item Le corps perd par rayonnement \res{\text{plus qu'il ne produit}} : la peau se refroidit et l'on a froid, même si l'air
est à \qty{20}{\degreeCelsius}. Le double vitrage, dont la face intérieure est plus chaude, améliore le confort.
\end{questions}
\end{corrige}
\end{blocLivre}
\begin{blocLivre}{cor-20-11}
\begin{corrige}{11}{Rafraîchir une canette}
\demarche{$C\,\dfrac{\mathrm{d}T}{\mathrm{d}t}=h\,S\,(T_\mathrm{ext}-T)$, puis $t=\tau\ln\dfrac{T_0-T_\mathrm{ext}}{T-T_\mathrm{ext}}$.}
\begin{questions}
\item Premier principe entre $t$ et $t+\mathrm{d}t$ : $C\,\mathrm{d}T=h\,S\,(T_\mathrm{ext}-T)\,\mathrm{d}t$, soit
$\res{\dfrac{\mathrm{d}T}{\mathrm{d}t}=-\dfrac{T-T_\mathrm{ext}}{\tau}}$ avec $\tau=\dfrac{C}{h\,S}$.
\item $C=0{,}330\times4180+13=\res{\qty{1,39e3}{J.K^{-1}}}$ et $\tau=\dfrac{1393}{10\times0{,}030}=\res{\qty{4,64e3}{s}}$, soit \qty{77}{min}.
\item $t=\tau\ln\dfrac{25-4}{8-4}=\res{\qty{2,1}{h}}$.
\item $t=\tau\ln\dfrac{25+18}{8+18}=\res{\qty{39}{min}}$. Oubliée, la canette atteindrait \qty{0}{\degreeCelsius} : l'eau, en gelant,
\res{\text{augmente de volume}} et peut la faire éclater.
\end{questions}
\end{corrige}
\end{blocLivre}
\begin{blocLivre}{cor-20-12}
\begin{corrige}{12}{Un modèle de l'effet de serre}
\demarche{$T_\mathrm{s}=T_\mathrm{e}\left(1-\dfrac{\varepsilon}{2}\right)^{-1/4}$.}
\begin{questions}
\item $T_\mathrm{s}=T_\mathrm{e}=\res{\qty{255}{K}}$ : sans absorption infrarouge, pas d'effet de serre.
\item $T_\mathrm{s}=2^{1/4}\,T_\mathrm{e}=\res{\qty{303}{K}}$ : trop chaud, l'atmosphère n'absorbe pas tout.
\item $1-\dfrac{\varepsilon}{2}=\left(\dfrac{254{,}6}{288}\right)^4$, d'où $\res{\varepsilon=0{,}78}$.
\item $T_\mathrm{s}$ passe de \qty{288,1}{K} à \qty{289,3}{K} : une hausse de \res{\qty{1,2}{K}} pour une faible variation de
$\varepsilon$. La température moyenne est très sensible à l'effet de serre.
\end{questions}
\end{corrige}
\end{blocLivre}
\begin{blocLivre}{cor-20-13}
\begin{corrige}{13}{Mars et Vénus}
\demarche{$S=\dfrac{1361}{d^2}$ ($d$ en ua) et $T=\left(\dfrac{(1-A)\,S}{4\,\sigma}\right)^{1/4}$.}
\begin{questions}
\item Mars : \res{\qty{586}{W.m^{-2}}} ; Vénus : \res{\qty{2,60e3}{W.m^{-2}}}.
\item Comme pour la Terre, $(1-A)\,S\,\pi R^2=\sigma\,T^4\times4\pi R^2$. Mars : \res{\qty{210}{K}} ; Vénus : \res{\qty{229}{K}}.
\item Pour Mars, la température mesurée est égale à la température d'équilibre : son atmosphère ténue ne produit
\res{\text{presque pas d'effet de serre}}. Vénus est plus chaude de \qty{508}{K} : son atmosphère épaisse de dioxyde de
carbone produit un \res{\text{effet de serre considérable}}.
\item Mercure n'a pas d'atmosphère ; Vénus, malgré un albédo élevé, \res{\text{piège}} le rayonnement infrarouge sous son
atmosphère dense.
\end{questions}
\end{corrige}
\end{blocLivre}
\begin{blocLivre}{cor-20-14}
\begin{corrige}{14}{La bouteille isotherme}
\demarche{$\dfrac{T_0-T_\mathrm{ext}}{T-T_\mathrm{ext}}=\mathrm{e}^{t/\tau}$, d'où $\tau=\dfrac{t}{\ln\dfrac{T_0-T_\mathrm{ext}}{T-T_\mathrm{ext}}}$.}
\begin{questions}
\item $\tau=\dfrac{\qty{6}{h}}{\ln\dfrac{70}{50}}=\res{\qty{17,8}{h}}$.
\item $\dfrac{17{,}8\times60}{25}\approx\res{43}$ fois plus lentement.
\item $60=20+70\,\mathrm{e}^{-t/\tau}$ donne $t=\tau\ln\dfrac{70}{40}=\res{\qty{10,0}{h}}$ ; à 12~h, $T=\res{\qty{55,7}{\degreeCelsius}}$.
\item $h\,S=\dfrac{C}{\tau}=\dfrac{\num{2.3e3}}{17{,}8\times3600}=\res{\qty{0,036}{W.K^{-1}}}$ : le vide supprime la conduction et la
convection entre les parois, l'argenture limite le rayonnement.
\end{questions}
\end{corrige}
\end{blocLivre}
\begin{blocLivre}{cor-20-15}
\begin{corrige}{15}{L'heure du crime}
\demarche{Rapport des écarts entre deux instants : $\mathrm{e}^{\Delta t/\tau}$ ; puis $t=\tau\ln\dfrac{37-18}{30-18}$ avant 23~h.}
\begin{questions}
\item $T(t)=18+(37-18)\,\mathrm{e}^{-t/\tau}$, d'où $\dfrac{T(t_1)-18}{T(t_2)-18}=\dfrac{\mathrm{e}^{-t_1/\tau}}{\mathrm{e}^{-t_2/\tau}}=\mathrm{e}^{(t_2-t_1)/\tau}$.
\item $\dfrac{12}{8}=\mathrm{e}^{2/\tau}$ ($t$ en heures), d'où $\tau=\dfrac{2}{\ln1{,}5}=\res{\qty{4,93}{h}}$.
\item Entre le décès et 23~h : $t=4{,}93\times\ln\dfrac{19}{12}=\qty{2,27}{h}$, soit 2~h~16~min : décès vers \res{\text{20~h~44}}.
\item Le refroidissement réel n'est pas exactement exponentiel (plateau initial, vêtements, corpulence, variations de la
température ambiante) : l'heure obtenue reste une \res{\text{estimation}}, entachée d'une incertitude importante.
\end{questions}
\end{corrige}
\end{blocLivre}
